% a06_BALtorus_ct7: Step 3(F), third part: the region-entropy estimate Ba2 (3.42) at block factor 5 (Sect. 3C of Ba2
% completed), with the counterterm moves of Ba3 p. 422 paid in it, and the g_k-localisation of Q^(k) (Ba2 p. 579).

\subsubsection*{The region sum \cite[(3.42), Sect.~3C]{Ba2} at block factor $5$}

\emph{Setting.} $d=3$, $L=5$, $M\ge1$, $s_b\defeq LM$ the side of a large block in lattice units of $T^{(k)}$ \cite[p.~607]{Ba1}; $R>R_0$, $\theta\defeq s_b/R\le\frac17$; $\varrho\ge1$ and $\bar p$ with $2\bar p>\varrho d+2$ (Ba{\l}aban's $\varrho=2$, $\bar p=6$: exponent $2\bar p-2-\varrho d=4$); $\eps_0\le e^{-1}$; $\hat l_k\defeq1+\ln t_k^{-1}$, $r_k\defeq r(t_k)=R\hat l_k^\varrho$, $p_k\defeq p(t_k)=b_0\hat l_k^{\bar p}$. The metric is the torus sup-metric \cite[(1.3)]{Ba1}; $N_w(X)$ is the closed $w$-neighbourhood; boxes are images of products of integer intervals. The torus $T^{(k)}=\Z^d/\Gamma_k$, $\Gamma_k=L\Gamma_{k+1}$, is aligned with the blocks: $\Gamma_k=\bigoplus_\mu P_{k,\mu}\Z$ with $s_b\mid P_{k,\mu}$ for $k\le K-1$ (a condition on the torus: its sides are multiples of the side of a large block at scale $K-1$, hence at every scale up to $K-1$; cf.\ \cref{def:s10:torus}, whose sides lie in $M\N$, and \cref{def:a07:btorus}); so the large blocks and the $L$-blocks tile $T^{(k)}$ and the prime $X\mapsto X'$ of \cite[(1.16)]{Ba1} is defined. Families and regions are as printed: at scale $k$ a family $F_k=(P_v,Q_v,R_v,P_s,Q_s,R_s)$ inside $\Lambda^{(k-1)'}_7$ ($\Lambda^{(-1)'}_7\defeq T^{(0)}$); each element $e$ has a point set $\mathrm{pts}(e)$ (a $P$-element its $L$-block; a bond or plaquette its ends or corners; an $R$-element its point) and a charge point $x_e\in\mathrm{pts}(e)$ (the lowest corner), with $\mathrm{pts}(e)\subset N_{L-1}(x_e)$; $S_k\defeq\bigcup_e\mathrm{pts}(e)$; $\Lambda^{(k)}_0$ is the union of the large blocks of $\Lambda^{(k-1)'}_7$ at distance $>r_k$ from $S_k$ \cite[p.~570]{Ba2}, $\Lambda^{(k)}_i$ by \cite[(2.8)]{Ba2}, ``admissible'' as on \cite[p.~566]{Ba2}. Write $X_k\defeq\Lambda^{(k)c}_0$, $\mathrm{Col}_k\defeq\Lambda^{(k-1)'}_7\cap\Lambda^{(k)c}_7$, $\mathrm{Sh}_k\defeq\Lambda^{(k-1)'}_5\cap\Lambda^{(k)c}_5$, $N_k\defeq|T^{(k)}|$, $I_k\defeq\{z\in T^{(k)}:\dist(z,T^{(k)}\setminus\Lambda^{(k-1)'}_7)\ge s_b\}$ ($I_0=T^{(0)}$). Constants: $c_{\mathrm{box}}\defeq\min\{a3^{-d}/8,\gamma_03^{-d}/4,3^{-d}\}$ (Ba{\l}aban's $c_0$ of \cite[p.~593]{Ba2}); $c_5,c_7\ge0$, $\kappa_0>0$ the $O(1)$'s and the exponent of \cite[(3.42)]{Ba2} (\cref{lem:a06f:C143}); $n_d\defeq4+2\max\{d,d(d-1)/2\}$ ($n_3=10$); $\nu_L\defeq14+(16+2L)\theta$, $\nu'_L\defeq\nu_L+2+3\theta$, $S_L\defeq\sum_{m\ge0}L^{-m}(1+m\ln L)^\varrho$, $C^0_D\defeq\nu'_L(S_L-1)$, $C_D\defeq C^0_D+14(1+\theta)$, $C_1\defeq n_d\ln2+d+1+c_5+c_7$, $\Theta_L\defeq C_1C_D^dR^d/(c_{\mathrm{box}}b_0^2\ln L)$, $C_{\mathrm{fin}}\defeq L^d(c_7\eps_0^{\kappa_0}+c_5)+e^{-(d+1)}\eps_0^{d+1}/(L-1)$. At $L=5$, $\varrho=2$, $\theta=\frac17$: $S_5=3.470$, $C^0_D=49.76$, $C_D=65.76$.

The printed \cite[Sect.~3C]{Ba2} is not followed: its partition \cite[(3.43)]{Ba2} (cubes with side in $(r,2r)$) need not exist at $L=5$, and the count \cite[(3.45)]{Ba2} with the first inequality of \cite[(3.46)]{Ba2} is false as printed: at $k=1$, with one $R_s$-element at a point $x_*$ of $T^{(0)}$ and $F_1=\emptyset$, $\Lambda^{(0)c}_7$ is a cube of large blocks whose prime is a union of scale-$1$ large blocks only if its corners are divisible by $L$, which fails for $L-1$ of every $L$ positions of $x_*$; then a scale-$1$ large block straddling the face lies in $\Lambda^{(1)c}_0$ and meets $\Lambda^{(0)'}_7$, so the printed $C_1$ is non-empty while $3^d(|P^{(1)}_v|+\dots)=0$, and the first inequality of \cite[(3.46)]{Ba2} reads $1\le\exp(-c_0p^2|C_1|)<1$. The count below is on the interior $I_k$ (\cref{lem:a06f:RC}), and the straddling layer is covered by the inherited boxes (\cref{lem:a06f:RD}).

\begin{lemma}[consequences of the printed rules]\label{lem:a06f:RE}
For every admissible sequence and $k\le K-1$: \textup{(a)} every point of every element of $F_k$ lies in $X_k$; \textup{(b)} $\Lambda^{(k)c}_i\subset\Lambda^{(k)c}_{i+1}$ and $\Lambda^{(k)c}_i\subset N_{i(r_k+s_b-1)}(X_k)$, $i=0,\dots,8$; \textup{(c)} if $z\in X_k\cap I_k$, the large block $\beta\ni z$ lies in $\Lambda^{(k-1)'}_7$ and $\dist(\beta,S_k)\le r_k$; \textup{(d)} $T^{(k)}\setminus\Lambda^{(k-1)'}_7=(\Lambda^{(k-1)c}_7)'\subset X_k$; \textup{(e)} $X_k,\mathrm{Col}_k,\mathrm{Sh}_k\subset\Lambda^{(k)c}_7$; \textup{(f)} if $X_i=\emptyset$ for all $i<j$ and $X_j\ne\emptyset$, then $\Lambda^{(j-1)'}_7=T^{(j)}$ and $X_j\cap I_j=X_j$.
\end{lemma}

\begin{proof}
(a) The large block $\beta$ of a point of $\mathrm{pts}(e)$ has $\dist(\beta,S_k)=0\le r_k$, so it is not one of the blocks of \cite[p.~570]{Ba2}; the large blocks tile $T^{(k)}$ and $\Lambda^{(k)}_0$ is a union of them. (b) \cite[(2.8)]{Ba2} adds the large blocks at distance $\le r_k$ (blocks of $\Lambda_i^c$ included); a point of such a block is within $(s_b-1)+r_k$ of $\Lambda^c_i$; induct. (c) $\beta$ has sup-diameter $s_b-1<s_b$, so $\beta\subset\Lambda^{(k-1)'}_7$, and being a candidate not chosen, $\dist(\beta,S_k)\le r_k$. (d) $\Lambda^{(k-1)}_7$ is a union of $L$-blocks, so $(X^c)'=T^{(k)}\setminus X'$; $\Lambda^{(k)}_0\subset\Lambda^{(k-1)'}_7$. (e) by (b). (f) $X_i=\emptyset$ gives $\Lambda^{(i)c}_7=\emptyset$ by (b), so $\Lambda^{(i)'}_7=T^{(i+1)}$.
\end{proof}

\begin{lemma}[number of families]\label{lem:a06f:RN}
For an admissible $\Lambda^{(k)}_0$ the number of families $F_k$ admissible for it is $\le2^{n_d|X_k|}$.
\end{lemma}

\begin{proof}
By \cref{lem:a06f:RE}(a) every element has all its points in $X_k$; charge each to its charge point; a point is the charge point of at most one $P_v$-, one $P_s$-, one $R_v$-, one $R_s$-element and of at most $\max\{d,d(d-1)/2\}$ bonds or plaquettes of each $Q$-type; a family is a subset of these $\le n_d|X_k|$ candidates.
\end{proof}

\begin{lemma}[the box partition]\label{lem:a06f:RP}
Let $q_k\defeq\lceil(r_k+s_b)/s_b\rceil$. For each $k\le K-1$ there is a partition $\mathcal Q_k$ of $T^{(k)}$ into boxes and a symmetric relation $Q\sim Q'$ (equal or adjacent) such that: (Q1) every box is a union of large blocks; (Q2) every box is the image of a product of $d$ integer intervals with at most $2r_k+3s_b\le(2+3\theta)r_k$ points each; (Q3a) every box has at most $3^d$ boxes $Q'\sim Q$; (Q3b) if $\dist(z,x)\le r_k+s_b$ then $Q(z)\sim Q(x)$.
\end{lemma}

\begin{proof}
In each direction $\mu$ let $m_\mu\defeq P_{k,\mu}/s_b$. If $m_\mu\ge q_k$, cut $\Z/P_{k,\mu}\Z$ into consecutive arcs, each a union of between $q_k$ and $2q_k-1$ consecutive large-block intervals $[js_b,(j+1)s_b)$ (write $m_\mu=nq_k+\rho'$, $0\le\rho'<q_k$, and give the $\rho'$ extra blocks to the $n$ arcs, at most $q_k-1$ each); if $m_\mu<q_k$ use one arc. The product of the lifted arc partitions is $\Gamma_k$-invariant, so its images form a partition of $T^{(k)}$; $Q\sim Q'$ if some lifts have arc indices differing by at most one in every coordinate. (Q1): arcs are unions of large-block intervals. (Q2): an arc has at most $(2q_k-1)s_b<2(r_k+s_b)+s_b$ points, or $P_{k,\mu}<q_ks_b\le r_k+2s_b$ in the one-arc case; $s_b\le\theta r_k$. (Q3a): the boxes $\sim Q$ are the images of the $3^d$ neighbours of one lift. (Q3b): take lifts with $|\tilde z-\tilde x|_\infty=\dist(z,x)\le r_k+s_b$; in a direction with $m_\mu\ge q_k$ every arc has $\ge q_ks_b\ge r_k+s_b$ points, so the two arcs are equal or consecutive; in a one-arc direction the indices agree.
\end{proof}

\begin{lemma}[the cube count on the interior]\label{lem:a06f:RC}
Put $\mathcal C_k\defeq\{Q\in\mathcal Q_k:Q\cap X_k\cap I_k\ne\emptyset\}$. For every family $F_k$ admissible for $\Lambda^{(k)}_0$, every $Q\in\mathcal C_k$ satisfies $Q\sim Q(x_e)$ for some element $e$ of $F_k$; hence $|\mathcal C_k|\le3^d|F_k|$, $|F_k|\defeq|P^{(k)}_v|+\dots+|R^{(k)}_s|$. At $k=0$ this is \cite[(3.44)]{Ba2}.
\end{lemma}

\begin{proof}
Let $z\in Q\cap X_k\cap I_k$; its large block $\beta$ lies in $Q$ (Q1); by \cref{lem:a06f:RE}(c) there are $z'\in\beta$, $y\in S_k$ with $\dist(z',y)\le r_k$, $y\in\mathrm{pts}(e)$ with $\dist(y,x_e)\le L-1<s_b$; so $\dist(z',x_e)\le r_k+s_b$ and $Q=Q(z')\sim Q(x_e)$ by (Q3b); count with (Q3a).
\end{proof}

\begin{lemma}[the bound on $\zeta''$]\label{lem:a06f:RZ}
$\zeta''_{\Lambda^{(k)}_0}\le2^{n_d|X_k|}\exp(-c_{\mathrm{box}}p_k^2|\mathcal C_k|)$.
\end{lemma}

\begin{proof}
The six coefficients $a/8,\gamma_0/4,\gamma_0/4,a/8,\gamma_0/4,1$ of \cite[(3.41)]{Ba2} are each $\ge3^dc_{\mathrm{box}}$, so each term of \cite[(3.41)]{Ba2} is $\le\exp(-3^dc_{\mathrm{box}}p_k^2|F_k|)\le\exp(-c_{\mathrm{box}}p_k^2|\mathcal C_k|)$ by \cref{lem:a06f:RC} ($\mathcal C_k$ depends on the region sequence only); the number of terms is \cref{lem:a06f:RN}.
\end{proof}

\begin{lemma}[weighted region entropy]\label{lem:a06f:RW}
For $\theta'\ge0$, $\sum_Xe^{-\theta'|X|}\le\exp(N_ke^{-\theta'})$, the sum over all unions $X$ of large blocks of $T^{(k)}$. With $\theta_k\defeq(d+1)\hat l_k$: $\sum_{k<K}N_ke^{-\theta_k}\le e^{-(d+1)}\eps_0^{d+1}|T^{(K)}|/(L-1)$.
\end{lemma}

\begin{proof}
The large blocks partition $T^{(k)}$ into $N_k/s_b^d$ blocks, so the sum is $(1+e^{-\theta's_b^d})^{N_k/s_b^d}\le\exp(N_ks_b^{-d}e^{-\theta's_b^d})\le\exp(N_ke^{-\theta'})$. Next $e^{-\theta_k}=e^{-(d+1)}t_k^{d+1}$ and $N_kt_k^d=|T^{(K)}|t_K^d$, so $\sum_{k<K}N_ke^{-\theta_k}=e^{-(d+1)}|T^{(K)}|t_K^d\sum_{k<K}t_k$ and $\sum_{k<K}t_k=t_K/(L-1)\le\eps_0/(L-1)$. (A $k$-independent price would give $\sum_kN_ke^{-\theta'}\sim\eps^{-d}|T|$: the logarithm in \cite[(3.47)]{Ba2} is exactly the price $\hat l_k$.)
\end{proof}

\begin{proposition}[from the sum to the supremum]\label{prop:a06f:RS}
Let $\mathcal F(\Lambda)$ be the summand of \cite[(3.42)]{Ba2} for the admissible sequence $\Lambda=(\Lambda^{(0)}_0,\dots,\Lambda^{(K-1)}_0)$ and $\Xi(\Lambda)\defeq\sum_{k<K}[(n_d\ln2+\theta_k)|X_k|+c_7t_k^{\kappa_0}|\mathrm{Col}_k|+c_5|\mathrm{Sh}_k|]$. For every set $\mathcal A$ of admissible sequences, $\sum_{\Lambda\in\mathcal A}\mathcal F(\Lambda)\le\exp(C_{\mathrm{fin}}|T^{(K)}|)\sup_{\Lambda\in\mathcal A}\exp(-\sum_{k<K}c_{\mathrm{box}}p_k^2|\mathcal C_k|+\Xi(\Lambda))$.
\end{proposition}

\begin{proof}
The $k=K$ terms of the two exponents of \cite[(3.42)]{Ba2} are $\le(c_7t_K^{\kappa_0}+c_5)L^d|T^{(K)}|$. By \cref{lem:a06f:RZ}, $\mathcal F(\Lambda)\le e^{L^d(\dots)|T^{(K)}|}\Psi(\Lambda)e^{-\sum_{k<K}\theta_k|X_k|}$ with $\Psi(\Lambda)\defeq\exp(-\sum_{k<K}c_{\mathrm{box}}p_k^2|\mathcal C_k|+\Xi(\Lambda))$; bound $\Psi$ by its supremum; $\Lambda\mapsto(X_0,\dots,X_{K-1})$ is injective and every $X_k$ is a union of large blocks, so $\sum_{\Lambda\in\mathcal A}e^{-\sum\theta_k|X_k|}\le\prod_{k<K}\sum_Xe^{-\theta_k|X|}\le\exp(\sum_{k<K}N_ke^{-\theta_k})$ by \cref{lem:a06f:RW}; collect into $C_{\mathrm{fin}}$.
\end{proof}

\begin{lemma}[{covering families; \cite[(3.48)--(3.52)]{Ba2} repaired}]\label{lem:a06f:RD}
Set $\mathcal D_0\defeq\mathcal C_0$ and, for $1\le k\le K-1$, $\mathcal D_k\defeq\{P_k(Q):Q\in\mathcal D_{k-1}\}\cup\mathcal C_k$, $P_k(Q)$ the prime of the union of the $L$-blocks of $T^{(k-1)}$ meeting $N_{w_{k-1}}(Q)$, $w_{k-1}\defeq7(r_{k-1}+s_b)+Ls_b$ (constructed on lifts; a box). Then for $k\le K-1$: \textup{(i)} $X_k\subset\bigcup\mathcal D_k$; \textup{(ii)} $|\mathcal D_k|\le\sum_{j\le k}|\mathcal C_j|$; \textup{(iii)} every box of $\mathcal D_k$ has side $\le C^0_Dr_k$; \textup{(iv)} $|\Lambda^{(k)c}_i|\le(C_Dr_k)^d\sum_{j\le k}|\mathcal C_j|$ for $i\le7$, and $|\Lambda^{(k)c}_8|\le(C^{(8)}_Dr_k)^d\sum_{j\le k}|\mathcal C_j|$ with $C^{(8)}_D\defeq C^0_D+16(1+\theta)$.
\end{lemma}

\begin{proof}
(ii) is immediate. (i) Induction on $k$; at $k=0$, $I_0=T^{(0)}$. Let $z\in X_k$. If $z\in I_k$, $z\in Q(z)\in\mathcal C_k$. Otherwise $\dist(z,T^{(k)}\setminus\Lambda^{(k-1)'}_7)\le s_b-1$, so by \cref{lem:a06f:RE}(d) there is $z_1$ with $B(z_1)\subset\Lambda^{(k-1)c}_7$ and $\dist(z,z_1)\le s_b-1$; by \cref{lem:a06f:RE}(b) at scale $k-1$ and the induction hypothesis every point $u$ of $B(z_1)$ lies in $N_{7(r_{k-1}+s_b)}(Q)$ for some $Q\in\mathcal D_{k-1}$; lifting, a point of $B(z)$ is within $L(s_b-1)$ of $u$, so $B(z)$ meets $N_{w_{k-1}}(Q)$ and $z\in P_k(Q)$. (iii) A box born at scale $j$ has side $s_j\le(2+3\theta)r_j$ (Q2); the union of the $L$-blocks meeting an interval of $n$ points has at most $n+2(L-1)$ points, so $s_{i+1}\le(s_i+2w_i+2L-2)/L$, and $2w_i+2L-2\le14r_i+(16+2L)s_b\le\nu_Lr_i$; iterating from $j<k$, $s_k\le L^{-(k-j)}s_j+\nu_L\sum_{i=j}^{k-1}L^{-(k-i)}r_i\le\nu'_L\sum_{i=j}^{k-1}L^{-(k-i)}r_i$, and $r_i\le(1+(k-i)\ln L)^\varrho r_k$ gives $s_k\le\nu'_L(S_L-1)r_k$; at $k=j$, $2+3\theta\le C^0_D$. (iv) By \cref{lem:a06f:RE}(b) and (i), $\Lambda^{(k)c}_i\subset\bigcup_{Q\in\mathcal D_k}N_{i(r_k+s_b)}(Q)$, each of at most $(C^0_Dr_k+2i(r_k+s_b))^d$ points; use (ii), with $i\le7$, resp.\ $i=8$.
\end{proof}

\begin{theorem}[{\cite[(3.42)]{Ba2} at $L=5$}]\label{thm:a06f:R7}
Let $c_5$, $c_7$, $\kappa_0$ be as in \cref{lem:a06f:C143}, let $c_8\ge0$ be free of $k$, $\eps$, $K$, the torus and the history, and add to the left side of \cite[(3.42)]{Ba2} the factor $\exp(\sum_{k<K}c_8|\Lambda^{(k)c}_8|)$ (the counterterm moves, \cref{prop:a06f:D2C}). Put $C_1^{(8)}\defeq C_1+c_8$, $\Theta^{(8)}_L\defeq C^{(8)}_1(C^{(8)}_D)^dR^d/(c_{\mathrm{box}}b_0^2\ln L)$ and assume
\[
(\mathrm{E.ent})\qquad(1+\ln(L/\eps_0))^{2\bar p-2-\varrho d}\ \ge\ \Theta^{(8)}_L.
\]
\textup{(a)} The (augmented) left side of \cite[(3.42)]{Ba2} is $\le\exp(C_{\mathrm{fin}}|T^{(K)}|)$. \textup{(b)} If (E.ent) holds with $2\Theta^{(8)}_L$, the part of the sum over the admissible sequences with some $X_k\ne\emptyset$ is $\le\exp(-c_{\mathrm{box}}p(t_{K-1})^2/2)\exp(C_{\mathrm{fin}}|T^{(K)}|)$. \textup{(c)} Both bounds are uniform in $\eps$, $K$, the torus and the history; (E.ent) is a cap on $\eps_0$ fixed after $d,a,L,M,R,b_0$ and after $c_5,c_7,c_8$, in Ba{\l}aban's order of choices. With $c_8=0$ this is \cite[(3.42)]{Ba2} as printed.
\end{theorem}

\begin{proof}
(a) By \cref{prop:a06f:RS} (with $\Xi$ augmented by $\sum_{k<K}c_8|\Lambda^{(k)c}_8|$, which is carried through its proof unchanged) it suffices that $\Xi(\Lambda)\le\sum_{j<K}c_{\mathrm{box}}p_j^2|\mathcal C_j|$. By \cref{lem:a06f:RE}(e), $t_k^{\kappa_0}\le1\le\hat l_k$ and $\Lambda^{(k)c}_7\subset\Lambda^{(k)c}_8$, every coefficient in $\Xi$ is $\le C^{(8)}_1\hat l_k$, so by \cref{lem:a06f:RD}(iv) and an exchange of sums, $\Xi\le C^{(8)}_1\sum_{k<K}\hat l_k|\Lambda^{(k)c}_8|\le C^{(8)}_1(C^{(8)}_D)^dR^d\sum_{j<K}|\mathcal C_j|\sum_{k=j}^{K-1}\hat l_k^{1+\varrho d}$. $\hat l_k$ decreases in $k$ and $K-j\le\log_L(\eps_0/t_j)\le\hat l_j/\ln L$, so the inner sum is $\le\hat l_j^{2+\varrho d}/\ln L$ and $\Xi\le c_{\mathrm{box}}b_0^2\Theta^{(8)}_L\sum_{j<K}\hat l_j^{2+\varrho d}|\mathcal C_j|$. For $j\le K-1$, $t_j\le\eps_0/L$, so $\hat l_j\ge1+\ln(L/\eps_0)$ and (E.ent) gives $\Theta^{(8)}_L\le\hat l_j^{2\bar p-2-\varrho d}$; hence $\Xi\le\sum_jc_{\mathrm{box}}b_0^2\hat l_j^{2\bar p}|\mathcal C_j|=\sum_jc_{\mathrm{box}}p_j^2|\mathcal C_j|$. (b) With $2\Theta^{(8)}_L$ the same chain gives $\Xi\le\frac12\sum_jc_{\mathrm{box}}p_j^2|\mathcal C_j|$; for a non-trivial sequence let $j$ be the first index with $X_j\ne\emptyset$; by \cref{lem:a06f:RE}(f), $X_j\cap I_j=X_j$, so $|\mathcal C_j|\ge1$, and $p_j\ge p_{K-1}$. (c) No step used a bound on $L$, $\eps$, $K$, the torus or the history beyond the setting.
\end{proof}

\begin{proposition}[the counterterm moves in the ledger]\label{prop:a06f:D2C}
Add at each step $k\le K-1$ the bounds of \cref{lem:a06d:D22,lem:a06d:D23} to the exponent that \cite[(2.43), (2.115)]{Ba2} carries into \cite[(2.117), (3.42)]{Ba2}. Then the left side of \cite[(3.42)]{Ba2} is bounded by \cref{thm:a06f:R7}(a), and its non-trivial part by \textup{(b)}, with $c_8\defeq Q_1(C_{\mathrm{D22}}+C_{\mathrm{D23}})$, $C_{\mathrm{D22}}$, $C_{\mathrm{D23}}$ the constants of those lemmas and $Q_1$ absorbing the factor $q(t_k)$ of a $\dm_{(4)}$ vertex (\cref{lem:a06e:T1}), and the $t^\kappa$ parts in the bulk sum \cite[(2.118)]{Ba2}: $C_\kappa\sum_kt_k^\kappa|T^{(k)}|\le C_\kappa\eps_0^{\kappa-d}(1-L^{d-\kappa})^{-1}|T|$ for $\kappa>d$.
\end{proposition}

\begin{proof}
$\Lambda^{(k)}_7\setminus\Lambda^{(k)}_8\subset\Lambda^{(k)c}_8$ and $T^{(k)}_1\setminus\Lambda^{(k)}_8=\Lambda^{(k)c}_8$, and $t_k^{1/2-\beta}\le1$, $t_kq(t_k)\le Q_1$, so the moves are $\le c_8|\Lambda^{(k)c}_8|+C_\kappa t_k^\kappa|T^{(k)}|$: the first part is a term of the augmented $\Xi$, the second a bulk term; they enter \cite[(2.117)]{Ba2} with the same sign and in the same place as the collar terms, and \cref{lem:a06f:C141} carries them through \cite[Sect.~3A]{Ba2}.
\end{proof}

\subsubsection*{The $g_k$-localisation of $Q^{(k)}$ \cite[p.~579]{Ba2}}

\begin{proposition}[the terms of $Q^{(k)}$ removed by $g_k$]\label{prop:a06f:H}
Let $1\le k\le K-1$, $\Omega_c\defeq B^k(\Lambda^{(k)}_0)$, $\Phi\defeq\theta_{k+1}A'^{(k)}$, and $Q^{(k)}[Y]\defeq\sum_{j<k}\sum_{n=1}^{\bar n}\frac1{n!}\frac{d^n}{d\tau^n}F_j(\Omega_c;\tilde B+\tau Y)|_{\tau=0}$ \cite[(2.105)--(2.106)]{Ba2}, in the setting of \cref{thm:a06d:G}. Let $c_g\defeq\sup|\partial^\eta g_k|$ ($\supp g_k\subset\Lambda^{(k)}_7$, $g_k=1$ on $\Lambda^{(k)}_8$). Then $|Q^{(k)}[\Phi]-Q^{(k)}[g_k\Phi]|\le F_H(t_k)(2M+1)^3|\Lambda^{(k)c}_8|$ with $F_H(t)\defeq C_Htp(t)^2(1+c_g)^2$, bounded on $(0,1]$ and free of $\eps$, $k$, $\eps_0$, the torus and the history. The other two replacements of \cite[(2.107)]{Ba2} (the background $\tilde B\to B^{(k+1),\eta}$ and the domain $\Omega_c\to B^k(\Lambda^{(k)}_2)$) are $O(t^\kappa)|\Lambda^{(k)}_7|$ for every $\kappa$.
\end{proposition}

\begin{proof}
Fix $j$ and $y$; put $\Phi_1\defeq g_k\Phi$, $\Phi_2\defeq(1-g_k)\Phi$, $\supp\Phi_2\subset\supp\theta_{k+1}\setminus B^k(\Lambda^{(k)}_8)$. For each $n$, $D^nQ_y(\tilde B)[\Phi^n]-D^nQ_y(\tilde B)[\Phi_1^n]$ is a sum of $n$ terms with at least one slot $\Phi_2$ (multilinearity). Gauge-fix at $y$ as in \cref{lem:a06d:G2} ($Q_y$ is gauge invariant, \cref{lem:a06d:G1}(b), and the axial map is linear). For $n=1$ use \cref{lem:a06d:G1}(f), which gives $D^2Q_y(s\beta)[\beta,\alpha_{\Phi_2}]$, bilinear in the curvature amplitudes of $\tilde B$ and $\Phi_2$; for $n\ge2$ use \cref{lem:a06d:G1}(e): each term is bounded by products of $n$ local amplitudes, one of them that of $\Phi_2$, plus terms suppressed by $e^{-\theta_2R/(4n)}$ with $R\defeq R_{kj}$ as in \cref{thm:a06d:G}. Curvatures in scale-$k$ units: $|\partial\Phi|\le Cp$ (the derivative clause of \cite[(2.91)]{Ba2}, Step~3(C$_1$): $\partial\Phi=(\partial\theta_{k+1})A'^{(k)}+\theta_{k+1}\partial A'^{(k)}$ with $|A'^{(k)}|\le C_{2.91}p$ by \cref{lem:a06d:C62}, $|\partial\theta_{k+1}|\le C$, and $\partial A'^{(k)}=\partial A^{(k)}-\partial B^{(k+1)}-\partial\tilde B^{(k)}$ by \cite[(2.89)]{Ba2}, bounded by (R1), \cite[(2.98)]{Ba2} and (R2) of Step~3(C$_2$)), $|\partial\Phi_2|\le|\partial\Phi|+2c_g|\Phi|\le C(1+c_g)p$, and on $N_k$, $\tilde B=A^{(k)}-\Phi-\tilde B^{(k)}$ (\cref{lem:a06d:G3}(a)) so $|\partial\tilde B|\le Cp$ by (R1), (R2); at scale $j$ all three amplitudes are $\le C(1+c_g)L^{-2(k-j)}t_k^{1/2}p$. Per $y$ the main terms are $\le C(1+c_g)^2L^{-4(k-j)}t_kp^2$, and the remaining terms are as in the proof of \cref{thm:a06d:G}; the $y$ within $R$ of $\supp\Phi_2$ number $\le(2R+3)^3(2M+1)^3L^{3(k-j)}|\Lambda^{(k)c}_8|$; sum over $y$ and over $j<k$ as there. The other two replacements act at distance $\ge5L^{k-j}r_k$ (scale-$j$ units) from $\supp g_k\Phi$, and \cref{lem:a06d:G1}(e) with one slot carrying the background or domain change gives $\le Ce^{-\theta_2L^{k-j}r_k/(2\bar n)}$ per $y$ (for the domain change, compare $F_j$ on the two domains through the common sub-box; the difference is localised at the domain boundary by \cite[(2.29)]{Ba1} and \cref{lem:a06d:G1}(d)).
\end{proof}
